\documentclass{article}

% packages
\usepackage{amsmath, amsthm, thmtools, amsfonts, amssymb, luacode, catchfile, hyperref, ifthen, graphicx, bbm}
\ifcsname c@kobocompile\endcsname
	\usepackage[a5paper, total={1072pt, 1448pt}, margin=10pt, includeheadfoot]{geometry} % set page margins
\else
	\usepackage[a4paper, margin=50pt, includeheadfoot]{geometry}
\fi
\usepackage[shortlabels]{enumitem}
\usepackage[skip=3pt, indent=0pt]{parskip}

% language
\usepackage[bidi=basic, layout=tabular, provide=*]{babel}
\ifcsname c@english\endcsname
	\babelprovide[main, import]{english}
\else
	\babelprovide[main, import, alph=letters.gershayim]{hebrew}
	\setlist[enumerate,2]{label={\alph*.}}
	\babelprovide{rl}
\fi
\babelfont{rm}[Renderer=Harfbuzz]{Libertinus Serif}
\babelfont{sf}{Libertinus Sans}
\babelfont{tt}{Libertinus Mono}

% style
\ifcsname c@stoppagebreaks\endcsname
\else
	\AddToHook{cmd/section/before}{\clearpage}	% Add line break before section
\fi
\linespread{1.3}
\setcounter{secnumdepth}{0}		% Remove default number tags from sections, this won't do well with theorems
\AtBeginDocument{\setlength{\belowdisplayskip}{3pt}}
\AtBeginDocument{\setlength{\abovedisplayskip}{3pt}}
\graphicspath{ {../images/} }

% operators
\DeclareMathOperator\cis{cis}
\DeclareMathOperator\Sp{Span}
\DeclareMathOperator\tr{tr}
\DeclareMathOperator\im{Im}
\DeclareMathOperator\re{Re}
\DeclareMathOperator\diag{diag}
\DeclareMathOperator*\lowlim{\underline{lim}}
\DeclareMathOperator*\uplim{\overline{lim}}
\DeclareMathOperator\rng{rng}
\DeclareMathOperator\Sym{Sym}
\DeclareMathOperator\Arg{Arg}
\DeclareMathOperator\Log{Log}
\DeclareMathOperator\dom{dom}
\DeclareMathOperator\supp{Supp}
\DeclareMathOperator\var{Var}
\DeclareMathOperator\cov{Cov}
\DeclareMathOperator\End{End}
\DeclareMathOperator\id{id}
\DeclareMathOperator\cl{cl}
\DeclareMathOperator\gal{Gal}
\DeclareMathOperator\dist{dist}
\DeclareMathOperator\form{form}
\DeclareMathOperator\term{term}
\DeclareMathOperator\sent{sent}
\DeclareMathOperator\Th{Th}
\DeclareMathOperator\Hom{Hom}
\DeclareMathOperator\Aut{Aut}
\DeclareMathOperator\Char{char}
\DeclareMathOperator\vol{vol}
\DeclareMathOperator\GL{GL}
\DeclareMathOperator\SL{SL}
\DeclareMathOperator\lcm{lcm}

% commands
\newcommand{\cupdot}{\mathbin{\mathaccent\cdot\cup}}
\newcommand{\indicator}{\mathbbm{1}}
\newcommand{\AAA}[0]{\mathbb{A}}
\newcommand{\Aa}[0]{\mathcal{A}}
\newcommand{\BB}[0]{\mathbb{B}}
\newcommand{\Bb}[0]{\mathcal{B}}
\newcommand{\CC}[0]{\mathbb{C}}
\newcommand{\Cc}[0]{\mathcal{C}}
\newcommand{\DD}[0]{\mathbb{D}}
\newcommand{\Dd}[0]{\mathcal{D}}
\newcommand{\EE}[0]{\mathbb{E}}
\newcommand{\FF}[0]{\mathbb{F}}
\newcommand{\Ff}[0]{\mathcal{F}}
\newcommand{\Ii}[0]{\mathcal{I}}
\newcommand{\Gg}[0]{\mathcal{G}}
\newcommand{\HH}[0]{\mathbb{H}}
\newcommand{\Hh}[0]{\mathcal{H}}
\newcommand{\Kk}[0]{\mathcal{K}}
\newcommand{\Ll}[0]{\mathcal{L}}
\newcommand{\Mm}[0]{\mathcal{M}}
\newcommand{\NN}[0]{\mathbb{N}}
\newcommand{\Nn}[0]{\mathcal{N}}
\newcommand{\PP}[0]{\mathbb{P}}
\newcommand{\Pp}[0]{\mathcal{P}}
\newcommand{\QQ}[0]{\mathbb{Q}}
\newcommand{\RR}[0]{\mathbb{R}}
\newcommand{\Rr}[0]{\mathcal{R}}
\newcommand{\Ss}[0]{\mathcal{S}}
\newcommand{\TT}[0]{\mathbb{T}}
\newcommand{\Tt}[0]{\mathcal{T}}
\newcommand{\UU}[0]{\mathbb{U}}
\newcommand{\Uu}[0]{\mathcal{U}}
\newcommand{\Vv}[0]{\mathcal{V}}
\newcommand{\Ww}[0]{\mathcal{W}}
\newcommand{\Xx}[0]{\mathcal{X}}
\newcommand{\Yy}[0]{\mathcal{Y}}
\newcommand{\ZZ}[0]{\mathbb{Z}}
\newcommand{\Zz}[0]{\mathcal{Z}}
\newcommand{\acts}[0]{\circlearrowright}
\newcommand{\explain}[2] {
	\begin{flalign*}
		 && \text{#2} && \text{#1}
	\end{flalign*}
}
\newcommand{\maketitleprint}[1][green]{
	\begin{center}
		\includegraphics[width=6cm]{dragon_#1}
	\end{center}
}

% theorem commands
\newtheoremstyle{c_remark}
	{}	% Space above
	{}	% Space below
	{}% Body font
	{}	% Indent amount
	{\bfseries}	% Theorem head font
	{}	% Punctuation after theorem head
	{.5em}	% Space after theorem head
	{\thmname{#1}\thmnumber{ #2}\thmnote{ \normalfont{\text{(#3)}}}}	% head content
\newtheoremstyle{c_definition}
	{3pt}	% Space above
	{3pt}	% Space below
	{}% Body font
	{}	% Indent amount
	{\bfseries}	% Theorem head font
	{}	% Punctuation after theorem head
	{.5em}	% Space after theorem head
	{\thmname{#1}\thmnumber{ #2}\thmnote{ \normalfont{\text{(#3)}}}}	% head content
\newtheoremstyle{c_plain}
	{3pt}	% Space above
	{3pt}	% Space below
	{\itshape}% Body font
	{}	% Indent amount
	{\bfseries}	% Theorem head font
	{}	% Punctuation after theorem head
	{.5em}	% Space after theorem head
	{\thmname{#1}\thmnumber{ #2}\thmnote{ \text{(#3)}}}	% head content

\ifcsname c@english\endcsname
	\theoremstyle{plain}
	\newtheorem{theorem}{Theorem}[section]
	\newtheorem*{theorem*}{Theorem}
	\newtheorem{lemma}[theorem]{Lemma}
	\newtheorem{proposition}[theorem]{Proposition}
	\newtheorem*{proposition*}{Proposition}

	\theoremstyle{definition}
	\newtheorem{definition}[theorem]{Definition}
	\newtheorem*{definition*}{Definition}
	\newtheorem{example}{Example}[section]
	\newtheorem{exercise}{Exercise}[section]

	\theoremstyle{remark}
	\newtheorem*{remark}{Remark}
	\newtheorem*{solution}{Solution}
	\newtheorem{corollary}[theorem]{Corollary}
	\newtheorem{notation}[theorem]{Notation}
\else
	\theoremstyle{c_plain}
	\newtheorem{theorem}{משפט}[section]
	\newtheorem*{theorem*}{משפט}
	\newtheorem{lemma}[theorem]{למה}
	\newtheorem{proposition}[theorem]{טענה}
	\newtheorem*{proposition*}{טענה}

	\theoremstyle{c_definition}
	\newtheorem{definition}[theorem]{הגדרה}
	\newtheorem*{definition*}{הגדרה}
	\newtheorem{example}{דוגמה}[section]
	\newtheorem{exercise}{תרגיל}[section]

	\theoremstyle{c_remark}
	\newtheorem*{remark}{הערה}
	\newtheorem*{solution}{פתרון}
	\newtheorem{corollary}[theorem]{מסקנה}
	\newtheorem{notation}[theorem]{סימון}
\fi

% Questions related commands
\newcounter{question}
\setcounter{question}{1}
\newcounter{sub_question}
\setcounter{sub_question}{1}

\ifcsname c@english\endcsname
	\newcommand{\question}[1][0]{
		\ifthenelse{#1 = 0}{}{\setcounter{question}{#1}}
		\section{Question \arabic{question}}
		\addtocounter{question}{1}
		\setcounter{sub_question}{1}
	}

	\newcommand{\subquestion}[1][0]{
		\ifthenelse{#1 = 0}{}{\setcounter{sub_question}{#1}}
		\subsection{Part \alph{sub_question}}
		\addtocounter{sub_question}{1}
	}
\else
	\newcommand{\question}[1][0]{
		\ifthenelse{#1 = 0}{}{\setcounter{question}{#1}}
		\section{שאלה \arabic{question}}
		\addtocounter{question}{1}
		\setcounter{sub_question}{1}
	}

	\newcommand{\subquestion}[1][0]{
		\ifthenelse{#1 = 0}{}{\setcounter{sub_question}{#1}}
		\subsection{סעיף \localecounter{letters.gershayim}{sub_question}}
		\addtocounter{sub_question}{1}
	}
\fi

% import lua and start of document
\directlua{common = require ('../common')}

\GetEnv{AUTHOR}

% headers
\author{\AUTHOR}
\date\today

\title{פתרון מטלה 01 --- יסודות בתורת ההצגות}

\DeclareMathOperator\reg{reg}

% chktex-file 44

\begin{document}
\maketitle
\maketitleprint[red]

\question{}
יהי $n \ge 2$ ונבחן את ההצגה ה־$n$־מימדית של $S_n$,
\[
	\rho : S_n \to \GL_n(\CC)
\]
המתקבלת על־ידי מיפוי של פרמוטציות $\sigma$ למטריצות פרמוטציה מתאימות.

\subquestion{}
נראה ש־$\rho$ מתפרקת לסכום של שתי הצגות $V_1 \oplus V_2$ כך ש־$\dim V_1 = 1$ וכן $\dim V_2 = n - 1$.
\begin{proof}
	כל מטריצת פרמוטציה $M$ מקיימת את התכונה שסכום כל שורה הוא $1$ וכן סכום כל עמודה הוא $1$, לכן זוהי מטריצה סטוכסטית ומקיימת את התכונה ש־$1$ ערך עצמי של המטריצה עבור הווקטור $v = {(1, \ldots, 1)}^t$.
	נגדיר את $V_1 = \Sp\{ v \}$ ולכן נקבל ש־$\rho(\sigma)(u) = u$ לכל $u \in V_1$ ולכל $\sigma \in S_n$.
	נגדיר את $V_2 = (e_1, \ldots, e_{n - 1})$ ולכן $V_1 \oplus V_2 = \CC^n$ בדיוק. \\
	מצאנו כבר שאם $u \in V_1$ אז גם $\sigma . u \in V_1$ ולכן נשאר להראות שאם $u \in V_2$ אז $\sigma . u \in V_2$ לכל $\sigma \in S_n$.
	נסמן $\pi_j : \CC^n \to \CC$ ההטלה על הציר ה־$j$, כלומר $\pi_j(e_i) = e_i \iff i = j$, ונבחין כי $\pi_n(u) = 0$, ולכן $\pi_{\sigma(n)}(u) = 0$ ובפרט $u \notin V_1$.
\end{proof}

\subquestion{}
נראה ש־$V_1$ ו־$V_2$ אי־פריקים.
\begin{proof}
	ברור כי $V_1$ אי־פריק, נראה זאת עבור $V_2$.
	נבחין כי לכל $1 \le i, j < n$ קיים $\sigma = (i\ j) \in S_n$ כך ש־$\sigma . e_i = e_j$, 
	יהי תת־מרחב $V \subset V_2$, לכן קיים $j < n$ כך ש־$e_j \notin U$, אחרת $\dim U = \dim V_2$ בסתירה להיותו תת־מרחב ממש, ויהי $0 \ne u \in U$, אז קיים אינדקס $i$ כך ש־$\pi_i(u) \ne 0$,
	לכן $\sigma . u \notin U$ עבור $\sigma$ שהגדרנו, ולכן $U$ לא מרחב אינווריאנטי, נסיק ש־$V_2$ עצמו הוא אי־פריק.
\end{proof}

\question{}
\subquestion{}
נוכיח שכל הצגה $(V, \rho)$ אי־פריקה של חבורה אבלית סופית $A$ מעל $\CC$ היא בהכרח ממימד $1$.
\begin{proof}
	נניח בשלילה ש־$\dim V \ge 2$ ולכן קיימים $u, v \in V \setminus \{ 0 \}$ בלתי־תלויים לינארית.
	ההצגה היא אי־פריקה ולכן קיים $\sigma \in A$ כך ש־$\sigma . u = v$, או לצירוף לינארי של הווקטורים ולכן נניח בלי הגבלת הכלליות את השוויון בלבד.
	נתון כי $A$ סופית ולכן קיים $n \in \NN$ כך ש־$\sigma^n = e$, בפרט $\sigma^n . u = u$ ונסיק שההעתקה $\rho(\sigma)$ היא נילפוטנטית ולכן אלכסונית, אבל זו סתירה ל־$\sigma . u = v$ שכן הווקטורים יהיו תלויים לינארית.
	זאת אפשר לראות על־ידי הנחה בלי הגבלת הכלליות ש־$u, v$ נבחרו כך שההעתקה $\rho(\sigma)$ היא בצורת ז'ורדן.
\end{proof}

\subquestion{}
נראה שזה לא בהכרח המצב עבור שדות אחרים.
\begin{proof}
	נגדיר את ההצגה $(\RR^2, \rho)$ מעל $\RR$ של החבורה $C_4 = \ZZ / 4 \ZZ$ הצגת הסיבוב, כלומר $C_4 = \langle \sigma \rangle$ עבור,
	\[
		\rho(\sigma)
		= \begin{pmatrix}
			0 & -1 \\
			1 & 0
		\end{pmatrix}
	\]
	נבחין כי בהרצאה כבר ראינו כי העתקה זו היא אי־פריקה.
\end{proof}

\question{}
\begin{definition}[הצגה רגולרית]
	ההצגה הרגולרית (השמאלית) של חבורה סופית $G$ היא הצגה $|G|$־מימדית מעל $\CC$ המוגדרת על־ידי הפעולה של $G$ על עצמה על־יד כפל (משמאל).
	במקרה זה נסמן $(V_{\reg}, \rho_{\reg})$ עבור $V_{\reg}$ מרחב לינארי מעל $\CC$ עם $\dim V_{\reg} = |G|$ והבסיס ${\{ v_g \}}_{g \in G}$ כך שמתקיים,
	\[
		g . v_h
		= \rho_{\reg}(g) \cdot v_h
		= v_{g h}
	\]
\end{definition}

\subquestion{}
נמצא עותק של ההצגה החד־מימדית הטריוויאלית של $G$ ב־$V_{\reg}$.
\begin{solution}
	נבחין כי $G$ סופית ולכן לכל $g \in G$ בהכרח $g^n = e$ לאיזשהו $n$, כלומר ${(\rho(g))}^n = \id$, לכן $\rho(g)$ היא העתקה סטוכסטית ו־$u = {(1, \ldots, 1)}^t$ וקטור עצמי של $\rho(g)$.
	נגדיר את $U = \Sp\{ u \}$ ונראה זאת.
\end{solution}

\end{document}
